% Options for packages loaded elsewhere
\PassOptionsToPackage{unicode}{hyperref}
\PassOptionsToPackage{hyphens}{url}
\PassOptionsToPackage{dvipsnames,svgnames,x11names}{xcolor}
%
\documentclass[
  letterpaper,
  DIV=11,
  numbers=noendperiod]{scrartcl}

\usepackage{amsmath,amssymb}
\usepackage{iftex}
\ifPDFTeX
  \usepackage[T1]{fontenc}
  \usepackage[utf8]{inputenc}
  \usepackage{textcomp} % provide euro and other symbols
\else % if luatex or xetex
  \usepackage{unicode-math}
  \defaultfontfeatures{Scale=MatchLowercase}
  \defaultfontfeatures[\rmfamily]{Ligatures=TeX,Scale=1}
\fi
\usepackage{lmodern}
\ifPDFTeX\else  
    % xetex/luatex font selection
\fi
% Use upquote if available, for straight quotes in verbatim environments
\IfFileExists{upquote.sty}{\usepackage{upquote}}{}
\IfFileExists{microtype.sty}{% use microtype if available
  \usepackage[]{microtype}
  \UseMicrotypeSet[protrusion]{basicmath} % disable protrusion for tt fonts
}{}
\makeatletter
\@ifundefined{KOMAClassName}{% if non-KOMA class
  \IfFileExists{parskip.sty}{%
    \usepackage{parskip}
  }{% else
    \setlength{\parindent}{0pt}
    \setlength{\parskip}{6pt plus 2pt minus 1pt}}
}{% if KOMA class
  \KOMAoptions{parskip=half}}
\makeatother
\usepackage{xcolor}
\setlength{\emergencystretch}{3em} % prevent overfull lines
\setcounter{secnumdepth}{5}
% Make \paragraph and \subparagraph free-standing
\makeatletter
\ifx\paragraph\undefined\else
  \let\oldparagraph\paragraph
  \renewcommand{\paragraph}{
    \@ifstar
      \xxxParagraphStar
      \xxxParagraphNoStar
  }
  \newcommand{\xxxParagraphStar}[1]{\oldparagraph*{#1}\mbox{}}
  \newcommand{\xxxParagraphNoStar}[1]{\oldparagraph{#1}\mbox{}}
\fi
\ifx\subparagraph\undefined\else
  \let\oldsubparagraph\subparagraph
  \renewcommand{\subparagraph}{
    \@ifstar
      \xxxSubParagraphStar
      \xxxSubParagraphNoStar
  }
  \newcommand{\xxxSubParagraphStar}[1]{\oldsubparagraph*{#1}\mbox{}}
  \newcommand{\xxxSubParagraphNoStar}[1]{\oldsubparagraph{#1}\mbox{}}
\fi
\makeatother


\providecommand{\tightlist}{%
  \setlength{\itemsep}{0pt}\setlength{\parskip}{0pt}}\usepackage{longtable,booktabs,array}
\usepackage{calc} % for calculating minipage widths
% Correct order of tables after \paragraph or \subparagraph
\usepackage{etoolbox}
\makeatletter
\patchcmd\longtable{\par}{\if@noskipsec\mbox{}\fi\par}{}{}
\makeatother
% Allow footnotes in longtable head/foot
\IfFileExists{footnotehyper.sty}{\usepackage{footnotehyper}}{\usepackage{footnote}}
\makesavenoteenv{longtable}
\usepackage{graphicx}
\makeatletter
\newsavebox\pandoc@box
\newcommand*\pandocbounded[1]{% scales image to fit in text height/width
  \sbox\pandoc@box{#1}%
  \Gscale@div\@tempa{\textheight}{\dimexpr\ht\pandoc@box+\dp\pandoc@box\relax}%
  \Gscale@div\@tempb{\linewidth}{\wd\pandoc@box}%
  \ifdim\@tempb\p@<\@tempa\p@\let\@tempa\@tempb\fi% select the smaller of both
  \ifdim\@tempa\p@<\p@\scalebox{\@tempa}{\usebox\pandoc@box}%
  \else\usebox{\pandoc@box}%
  \fi%
}
% Set default figure placement to htbp
\def\fps@figure{htbp}
\makeatother
% definitions for citeproc citations
\NewDocumentCommand\citeproctext{}{}
\NewDocumentCommand\citeproc{mm}{%
  \begingroup\def\citeproctext{#2}\cite{#1}\endgroup}
\makeatletter
 % allow citations to break across lines
 \let\@cite@ofmt\@firstofone
 % avoid brackets around text for \cite:
 \def\@biblabel#1{}
 \def\@cite#1#2{{#1\if@tempswa , #2\fi}}
\makeatother
\newlength{\cslhangindent}
\setlength{\cslhangindent}{1.5em}
\newlength{\csllabelwidth}
\setlength{\csllabelwidth}{3em}
\newenvironment{CSLReferences}[2] % #1 hanging-indent, #2 entry-spacing
 {\begin{list}{}{%
  \setlength{\itemindent}{0pt}
  \setlength{\leftmargin}{0pt}
  \setlength{\parsep}{0pt}
  % turn on hanging indent if param 1 is 1
  \ifodd #1
   \setlength{\leftmargin}{\cslhangindent}
   \setlength{\itemindent}{-1\cslhangindent}
  \fi
  % set entry spacing
  \setlength{\itemsep}{#2\baselineskip}}}
 {\end{list}}
\usepackage{calc}
\newcommand{\CSLBlock}[1]{\hfill\break\parbox[t]{\linewidth}{\strut\ignorespaces#1\strut}}
\newcommand{\CSLLeftMargin}[1]{\parbox[t]{\csllabelwidth}{\strut#1\strut}}
\newcommand{\CSLRightInline}[1]{\parbox[t]{\linewidth - \csllabelwidth}{\strut#1\strut}}
\newcommand{\CSLIndent}[1]{\hspace{\cslhangindent}#1}

\KOMAoption{captions}{tableheading}
\makeatletter
\@ifpackageloaded{caption}{}{\usepackage{caption}}
\AtBeginDocument{%
\ifdefined\contentsname
  \renewcommand*\contentsname{Table of contents}
\else
  \newcommand\contentsname{Table of contents}
\fi
\ifdefined\listfigurename
  \renewcommand*\listfigurename{List of Figures}
\else
  \newcommand\listfigurename{List of Figures}
\fi
\ifdefined\listtablename
  \renewcommand*\listtablename{List of Tables}
\else
  \newcommand\listtablename{List of Tables}
\fi
\ifdefined\figurename
  \renewcommand*\figurename{Figure}
\else
  \newcommand\figurename{Figure}
\fi
\ifdefined\tablename
  \renewcommand*\tablename{Table}
\else
  \newcommand\tablename{Table}
\fi
}
\@ifpackageloaded{float}{}{\usepackage{float}}
\floatstyle{ruled}
\@ifundefined{c@chapter}{\newfloat{codelisting}{h}{lop}}{\newfloat{codelisting}{h}{lop}[chapter]}
\floatname{codelisting}{Listing}
\newcommand*\listoflistings{\listof{codelisting}{List of Listings}}
\makeatother
\makeatletter
\makeatother
\makeatletter
\@ifpackageloaded{caption}{}{\usepackage{caption}}
\@ifpackageloaded{subcaption}{}{\usepackage{subcaption}}
\makeatother

\usepackage{bookmark}

\IfFileExists{xurl.sty}{\usepackage{xurl}}{} % add URL line breaks if available
\urlstyle{same} % disable monospaced font for URLs
\hypersetup{
  pdftitle={Power Analysis of Subgroup Representation in Clinical Trials: Evidence from the ACTG 175 Trial},
  pdfauthor={Susan Peck, Joseph Burks, Jan Blackburn},
  colorlinks=true,
  linkcolor={blue},
  filecolor={Maroon},
  citecolor={Blue},
  urlcolor={Blue},
  pdfcreator={LaTeX via pandoc}}


\title{Power Analysis of Subgroup Representation in Clinical Trials:
Evidence from the ACTG 175 Trial}
\author{Susan Peck, Joseph Burks, Jan Blackburn}
\date{}

\begin{document}
\maketitle
\begin{abstract}
This study evaluates whether subgroup representation in clinical trials
provides sufficient statistical power to detect meaningful treatment
effects. Using the ACTG 175 trial as a case study, we examine how
subgroup size and event counts influence the ability to detect
heterogeneity in treatment response.
\end{abstract}


\section{Introduction}\label{introduction}

Randomized clinical trials are typically powered to detect treatment
effects in the overall study population rather than within demographic
subgroups. This can lead to underrepresentation for female and minority
ethnic groups, reducing the effective power of statistical tests and the
ability to detect critical differences in treatment effects. Prior work
has shown underrepresentation of women and ethnic minority groups in
pivotal clinical trials across therapeutic areas
(\citeproc{ref-khan2020women}{Khan et al. 2020};
\citeproc{ref-poon2013participation}{Poon et al. 2013};
\citeproc{ref-eshera2015demographics}{Eshera et al. 2015}).

A key feature of many clinical trials, including ACTG 320, is that
outcomes are measured as time-to-event outcomes and analyzed using
survival analysis methods. In this framework, statistical inference is
driven by the number of observed events (e.g., disease progression or
death) instead of the number of participants enrolled. As a result,
subgroups with smaller sample sizes tend to yield fewer events, which
can substantially reduce the power to detect treatment effects.

This distinction between sample size and event count is central to the
present study. Even when treatment effects differ across demographic
groups, studies may lack sufficient statistical power to detect these
differences if the number of observed events is small. This creates a
structural limitation in subgroup analysis that is not always apparent
when focusing solely on enrollment counts.

The purpose of this paper is to use the ACTG 320 clinical trial
(\citeproc{ref-hammer1997controlled}{Hammer et al. 1997}) as a case
study to investigate how underrepresentation in gender and ethnic
subgroups impacts the statistical power to detect clinically meaningful
treatment effects.

The rest of this paper discusses the ACTG320 data set elements and
limitations, provides a brief overview of key survival analysis concepts
including hazard functions, Cox proportional hazards models and how
hypothesis testing and power analysis are applied in this context. The
paper concludes with a discussion of the results, implications and areas
for further research.

\section{Data}\label{data}

The ACTG 320 dataset is derived from a randomized, double-blind clinical
trial evaluating the efficacy of combination antiretroviral therapy in
HIV-infected patients with advanced disease (CD4 cell counts ≤ 200
cells/mm³). Participants were stratified by their CD4 cell counts at
enrollment for those with CD4 cell count \textless{} 50 and patients
with CD4 cell count 51 - 200 cells/mm³). Patients in each group were
randomly assigned to receive either a two-drug nucleoside regimen with a
placebo or a three-drug regimen that included the protease inhibitor
indinavir.

The study's primary endpoint is a time-to-event defined as either
progression to AIDS or death. A total of 1,156 patients participated in
the trial but only 96 experienced progression to AIDS or death during
follow-up. Researchers also investigated two secondary endpoints:
changes in CD40 cell levels tracked at intervals during the study and
changes in plasma HIV-1 RNA concentrations.

The dataset for this paper was obtained from the GLDreg R library
(\citeproc{ref-GLDreg}{Su et al. 2023}). The file contains participant
level records that include treatment assignment, baseline clinical
characteristics (e.g., CD4 count, prior antiretroviral therapy),
demographic variables, and follow-up information on clinical. There are
also subgroup indicators (e.g., sex and race/ethnicity), which allow for
investigation of treatment effect heterogeneity and, importantly for
this study, assessment of how reduced subgroup sample sizes and event
counts impact statistical power.

The publicly available data set does not include data on the secondary
outcomes so this paper will focus on power analysis on primary endpoint
only. Considering that there 96 of enrolled participants experienced the
AIDS or death, the power of any subgroup tests is expected to be low,
limiting the precision of any subgroup level estimation treatment
effects.

\section{Methods}\label{methods}

Clinical trials evaluating time-to-event outcomes require statistical
methods that properly account for censoring and the timing of events. In
the context of HIV clinical trials such as ACTG 320, the primary
endpoint is the time to progression to AIDS or death, making survival
analysis the appropriate framework. The section will start with a review
of key concepts in survival analysis and progress to regression model
formulations, hypothesis testing and power analysis definitions that are
applicable.

\subsection{Survival and Hazard
Functions}\label{survival-and-hazard-functions}

Let \(T\) denote a non-negative random variable representing the time to
the event of interest. The survival function is defined as

\[
S(t) = P(T > t).
\]

This quantity represents the probability that an individual remains
event-free beyond time \(t\) and provides the basis for survival and
hazard analysis. The hazard function is defined as

\[
h(t) = \lim_{\Delta t \to 0} \frac{P(t \le T < t+\Delta t \mid T > t)}{\Delta t} \qquad (1).
\]

where \(\Delta t\) is a small time interval, and \(h(t)\) represents the
instantaneous event rate at time \(t\) conditional on survival up to
time \(t\) (\citeproc{ref-moore2016_ch2}{Moore 2016}). Unlike the
survival function, the hazard function is not a probability but
determines how quickly the survival probability declines over time.
Higher values of \(h(t)\) correspond to a greater risk of the event
occurring immediately after time \(t\), while lower values indicate
slower event occurrence. The hazard and survival functions are related
by

The hazard function can be derived from (1) as the ratio of the pdf of T
at time t and the survival function

\[
h(t) = \frac{f(t)}{S(t)}.
\]

\subsection{Cox Proportional Hazards Model and Hazard
Ratios}\label{cox-proportional-hazards-model-and-hazard-ratios}

The Cox proportional hazards model are used in time-to-event clinical
studies to estimate the association between covariates and the hazard of
the event. The model is semi-parametric since it doesn't require
specifying the baseline hazard function. For subject \(i\) with
covariate vector \({X}_i\), the Cox model is

\[
h_i(t \mid \mathbf{x}_i) = h_0(t) \exp(\mathbf{x}_i^\top \boldsymbol{\beta}),
\] where \({h}_0(t)\) is the unspecified baseline hazard function and
\(Beta\) is a vector of regression coefficients. In the Cox model,
comparisons between groups are made through ratios of hazard functions,
allowing treatment effects to be interpreted as relative differences in
instantaneous risk over time (\citeproc{ref-moore2016_ch2}{Moore 2016}).
The hazard function compares the hazard function between two groups

\[
HR = frac{h_T(t) / {h_C(t)}
\] where \(h_T(t)\) is the hazard function for the treatment group and
\(h_C(t)\) hazard function for the control group.

A hazard ratio less than one indicates the treatment reduces the
instantaneous risk versus the control and a value greater than one
indicates a higher instantaneous risk relative to control. Parameters
are estimated. Since the hazard ratios don't depend on \(t\), they are
assumed to be constant over time.

The Cox proportional hazards model is estimated using partial likelihood
estimation to account for censoring of the data
(\citeproc{ref-moore_ch5}{\textbf{moore\_ch5?}}). Let \(R(t_i)\) denote
the risk set immediately before event time \(t_i\). The partial
likelihood is

\[
L(\beta) = \prod_{i:\delta_i=1} 
\frac{\exp\left(X_i^\top \beta\right)}
{\sum_{j \in R(t_i)} \exp\left(X_j^\top \beta\right)}.
\]

The partial likelihood is converted to the log-likelihood which is
maximized to obtain the estimate \(\hat{\beta}\).

\subsection{Likelihood Ratio Test}\label{likelihood-ratio-test}

The likelihood ratio test is generally preferred in survival analysis,
particularly when the number of observed events is limited, as it relies
on comparisons of overall model fit rather than asymptotic normality of
individual parameter estimates
(\citeproc{ref-moore_ch5}{\textbf{moore\_ch5?}}). Let \(\ell(\beta)\)
denote the log partial likelihood of the Cox model. To test a hypothesis
about a subset of parameters, two models are considered:

\begin{itemize}
\tightlist
\item
  A \textbf{reduced model}, which imposes the null hypothesis
  constraints, and\\
\item
  A \textbf{full model}, which includes all parameters of interest.
\end{itemize}

The likelihood ratio test (LRT) statistic is defined as

\[
\Lambda = -2\left[\ell(\hat{\beta}_{\text{reduced}}) - \ell(\hat{\beta}_{\text{full}})\right].
\] Under standard regularity conditions, and assuming the null
hypothesis is true, the test statistic follows an approximate chi-square
distribution:

\[
\Lambda \sim \chi^2_k,
\]

where \(k\) is the number of restrictions imposed under the null
hypothesis. For example, \(k = 1\) for a single interaction term and
\(k > 1\) for joint tests involving multiple subgroup effects.

The null hypothesis is rejected at significance level \(\alpha\) if

\[
\Lambda > \chi^2_{k,\,1-\alpha}.
\]

\subsection{Power Estimation for Event-Based
Studies}\label{power-estimation-for-event-based-studies}

Power calculations for time-to-event analyses are driven primarily by
the number of observed events rather than the total sample size. For a
two-group comparison, the approximate number of events required to
detect a specified hazard ratio \(HR\) with significance level
\(\alpha\) and power \(1-\beta\) is

For a two-group survival comparison, the approximate power for a fixed
number of observed events is

\[
\text{Power}
=
\Phi\left(
\sqrt{D p(1-p)}\,|\log(HR)|
-
z_{1-\alpha/2}
\right)
\]

where \(D\) is the observed or planned number of events, \(p\) is the
proportion assigned to treatment, \(HR\) is the target hazard ratio,
\(\alpha\) is the two-sided significance level, and \(\Phi(\cdot)\) is
the standard normal cumulative distribution function.

For equal allocation, where \(p = 0.50\), this becomes

\[
\text{Power}
=
\Phi\left(
\frac{\sqrt{D}}{2}|\log(HR)|
-
z_{1-\alpha/2}
\right).
\]

This expression highlights that the ability to detect treatment effects
depends on both the magnitude of the hazard ratio and the total number
of events observed during follow-up. Smaller effect sizes require
substantially more events to achieve adequate power
(\citeproc{ref-Schoenfeld1983}{Schoenfeld 1983}).

\subsection{Primary Treatment Effect
Model}\label{primary-treatment-effect-model}

The overall treatment effect may be evaluated using a Cox proportional
hazards model of the form

\[
h_i(t) =
h_0(t)\exp\left(
\beta_1 T_i
+ \boldsymbol{\beta}_2^{\top}Z_i
\right)
\]

where \(T_i\) is the treatment indicator and \(Z_i\) is a vector of
additional baseline covariates. The coefficient \(\beta_1\) represents
the log hazard ratio for the primary treatment effect after adjusting
for the covariates in \(Z_i\). The treatment hazard ratio is

\[
HR_T = \exp(\beta_1)
\]

The null hypothesis for the overall treatment effect is

\[
\begin{gathered}
H_0 : \beta_1 = 0 \quad \text{vs.} \quad H_A : \beta_1 \neq 0 \\
\text{or equivalently} \\
H_0 : HR_T = 1 \quad \text{vs.} \quad H_A : HR_T \neq 1
\end{gathered}
\]

\subsection{Measuring Subgroup Treatment Effects Between
Genders}\label{measuring-subgroup-treatment-effects-between-genders}

To evaluate whether the treatment effect differs between males and
females, the regression model is expanded to include a gender indicator
and a treatment-by-gender interaction term:

\[
h_i(t) =
h_0(t)\exp\left(
\beta_1 T_i
+ \beta_2 G_i
+ \beta_3 T_iG_i
+ \boldsymbol{\beta}_4^{\top}Z_i
\right)
\]

where \(G_i\) is a binary gender indicator and \(Z_i\) represents
additional covariates, \(\beta_1\) is the treatment effect for the
reference gender group, \(\beta_2\) is the baseline difference in hazard
between gender groups among participants in the control treatment group,
and \(\beta_3\) is the treatment-by-gender interaction effect.

The treatment hazard ratio for the reference gender group is

\[
HR_{T \mid G = 0} = \exp(\beta_1)
\]

The treatment hazard ratio for the comparison gender group is

\[
HR_{T \mid G = 1} = \exp(\beta_1 + \beta_3)
\]

The hypothesis test for whether the treatment effect differs by gender
is

\[
\begin{gathered}
H_0 : \beta_3 = 0 \quad \text{vs.} \quad H_A : \beta_3 \neq 0 \\
\text{or equivalently} \\
H_0 : exp(\beta_3) = 1 \quad \text{vs.} \quad H_A : exp(\beta_3) \neq 1
\end{gathered}
\]

\subsection{Measuring Subgroup Treatment Effects Between Ethnic
Groups}\label{measuring-subgroup-treatment-effects-between-ethnic-groups}

The same interaction framework can be extended to evaluate whether
treatment effects differ across more than two ethnic groups. In the
ACTG320 trial, ethnicity is coded as follows: 1 = White non-Hispanic, 2
= Black non-Hispanic, 3 = Hispanic, 4 = Asian or Pacific Islander, 5 =
American Indian or Alaska Native, and 6 = Other or unknown.

Using White non-Hispanic participants as the reference group, the model
is

\[
\begin{aligned}
h_i(t) = h_0(t)\exp\big(&
\beta_1 T_i
+ \beta_2 Black_i
+ \beta_3 Hispanic_i
+ \beta_4 AsianPI_i
+ \beta_5 AIAN_i
+ \beta_6 Other_i \\
&+ \beta_7 T_iBlack_i
+ \beta_8 T_iHispanic_i
+ \beta_9 T_iAsianPI_i
+ \beta_{10} T_iAIAN_i
+ \beta_{11} T_iOther_i
+ \boldsymbol{\beta}_{12}^{\top}Z_i
\big)
\end{aligned}
\]

In this model, \(\beta_1\) is the treatment effect for White
non-Hispanic participants. The coefficients \(\beta_2\) through
\(\beta_6\) represent baseline hazard differences between each ethnic
subgroup and the White non-Hispanic reference group among participants
in the control treatment group.

The interaction coefficients \(\beta_7\) through \(\beta_{11}\) measure
whether the treatment effect differs between each ethnic subgroup and
the White non-Hispanic reference group. The vector \(Z_i\) includes
additional baseline covariates used for adjustment but not for subgroup
treatment-effect testing.

The joint hypothesis test for whether treatment effects differ across
ethnic groups is

\[
\begin{gathered}
H_0 : \beta_7 = \beta_8 = \beta_9 = \beta_{10} = \beta_{11} = 0 \\
H_A : \text{At least one interaction coefficient is not equal to zero}
\end{gathered}
\]

If the null hypothesis is rejected, there is evidence that the treatment
effect differs for at least one ethnic subgroup relative to the White
non-Hispanic reference group.

\subsection{Power Analysis for Treatment and Subgroup
Effects}\label{power-analysis-for-treatment-and-subgroup-effects}

Power calculations for survival outcomes are based on the number of
observed events rather than the total sample size. Under the Cox
proportional hazards model, the statistical power to detect treatment
effects is driven by the magnitude of the hazard ratio and the total
number of events observed during follow-up.

\subsection{Power for the Overall Treatment
Effect}\label{power-for-the-overall-treatment-effect}

For a two-group survival comparison, power can be estimated directly
when the number of observed or anticipated events is known. This
formulation is useful for sensitivity analysis because the available
event count may be fixed by the completed study or by the expected
duration of follow-up.

For a two-sided test at significance level \(\alpha\), the approximate
power is

\[
\text{Power}
=
\Phi\left(
\sqrt{D p(1-p)}\,|\log(HR)|
-
z_{1-\alpha/2}
\right),
\]

where \(D\) is the total number of observed or anticipated events, \(p\)
is the proportion assigned to the treatment group, \(1-p\) is the
proportion assigned to the control group, \(HR\) is the target hazard
ratio comparing treatment with control, \(z_{1-\alpha/2}\) is the
standard normal critical value for a two-sided test, and \(\Phi(\cdot)\)
is the standard normal cumulative distribution function.

For equal allocation, where \(p = 0.50\), this simplifies to

\[
\text{Power}
=
\Phi\left(
\frac{\sqrt{D}}{2}|\log(HR)|
-
z_{1-\alpha/2}
\right).
\] This expression estimates the probability of rejecting the null
hypothesis of no treatment effect,

\[
H_0: HR = 1,
\]

when the true treatment effect is assumed to be \(HR \neq 1\). This is
the inverse form of the standard event-count formula for survival
analysis and is appropriate for assessing whether the available number
of primary endpoint events provides adequate power for detecting a
specified hazard ratio..

\subsection{Power for Subgroup Differences:
Gender}\label{power-for-subgroup-differences-gender}

For subgroup interaction analyses, power can be estimated directly from
the anticipated number of observed events. Because interaction effects
are estimated from contrasts between subgroup-specific treatment
effects, the effective amount of information is reduced relative to the
overall treatment comparison. Consequently, substantially larger event
counts are typically required to detect treatment effect heterogeneity
across subgroups (\citeproc{ref-Hsieh2000}{Hsieh and Lavori 2000}).

For a two-sided interaction test at significance level \(\alpha\), the
approximate power is

\[
\text{Power}
=
\Phi\left(
\sqrt{
D \,
p_T(1-p_T)\,
p_G(1-p_G)
}
\,
|\log(HR_{\text{int}})|
-
z_{1-\alpha/2}
\right),
\]

where \(D\) is the total number of observed or anticipated events,
\(p_T\) is the proportion assigned to the treatment group, \(p_G\) is
the proportion of participants in one of the subgroup categories (for
example, one gender subgroup), \(HR_{\text{int}}\) is the interaction
hazard ratio, \(z_{1-\alpha/2}\) is the standard normal critical value
for a two-sided test, and \(\Phi(\cdot)\) is the standard normal
cumulative distribution function (\citeproc{ref-Hsieh2000}{Hsieh and
Lavori 2000})..

The interaction hazard ratio quantifies the relative difference between
subgroup-specific treatment effects. For example, if

\[
HR_1
\]

denotes the treatment hazard ratio in subgroup 1 and

\[
HR_0
\]

denotes the treatment hazard ratio in subgroup 0, then the interaction
hazard ratio is

\[
HR_{\text{int}}
=
\frac{HR_1}{HR_0}.
\]

Equivalently,

\[
\log(HR_{\text{int}})
=
\log(HR_1)
-
\log(HR_0).
\]

For equal treatment allocation and equally sized subgroups, where

\[
p_T = 0.50
\quad \text{and} \quad
p_G = 0.50,
\]

the power approximation simplifies to

\[
\text{Power}
=
\Phi\left(
\frac{\sqrt{D}}{4}
|\log(HR_{\text{int}})|
-
z_{1-\alpha/2}
\right).
\]

This formulation estimates the probability of rejecting the null
hypothesis of no treatment-by-subgroup interaction,

\[
H_0: HR_{\text{int}} = 1,
\]

when the true interaction effect satisfies

\[
HR_{\text{int}} \neq 1.
\]

\subsection{Subgroup Power Analysis:
Ethnicity}\label{subgroup-power-analysis-ethnicity}

For ethnic subgroup interaction analyses, power can be estimated
directly from the anticipated number of observed events. Similar to
gender subgroup analyses, interaction effects are estimated from
contrasts between subgroup-specific treatment effects, reducing the
effective amount of information available relative to the overall
treatment comparison. Because multiple ethnic subgroup contrasts are
evaluated simultaneously, substantially larger event counts are
generally required to detect meaningful treatment effect heterogeneity
across ethnic groups (\citeproc{ref-Hsieh2000}{Hsieh and Lavori 2000}).

For a two-sided interaction test at significance level \(\alpha\), the
approximate power is

\[
\text{Power}
=
\Phi\left(
\sqrt{
D \,
p_T(1-p_T)\,
p_E(1-p_E)
}
\,
|\log(HR_{\text{int}})|
-
z_{1-\alpha/2}
\right),
\] \(p_T\) is the proportion assigned to the treatment group, \(p_E\) is
the proportion of participants in one of the ethnic subgroup categories,
\(HR_{\text{int}}\) is the interaction hazard ratio comparing
subgroup-specific treatment effects, \(z_{1-\alpha/2}\) is the standard
normal critical value for a two-sided test, and \(\Phi(\cdot)\) is the
standard normal cumulative distribution function
(\citeproc{ref-Hsieh2000}{Hsieh and Lavori 2000}).

\subsection{Power Analysis Simulation}\label{power-analysis-simulation}

Since there was no power analysis associated with this study, a baseline
power analysis on overall treatment effect will be performed as point of
reference for power analysis by gender and ethnic subgroups.

To gain a better understanding of number of patients needed to be
recruited by subgroup to achieve desired levels of power a simulation
will be run to estimate number of events, and indirectly the number of
patients that would need to be recruited by subgroup to achieve the
desired power.

Since there are multiple hypotheses being tested on the same data set, a
Bonferroni adjustment will be applied to the significance level of each
test to achieve a family wise error rate of \(\alpha\) = 0.05. Since we
are testing three distinct hypotheses, this will be accomplished by
applying a Bonferroni adjustment to the significance used in each
individual test statistic so that \$\alpha\_adj = 0.05 / 3 = 0.017.

\section{Results}\label{results}

The analysis in this paper was completed using the R statistical package
(\citeproc{ref-R-base}{R Core Team 2026}) primarily using the survival
library (\citeproc{ref-survival-package}{Therneau 2026}).

\section{Conclusion}\label{conclusion}

Brilliant and compelling conclusion

\phantomsection\label{refs}
\begin{CSLReferences}{1}{0}
\bibitem[\citeproctext]{ref-eshera2015demographics}
Eshera, N., H. Itana, L. Zhang, G. Soon, and E. Fadiran. 2015.
{``Demographics of Clinical Trial Participants in Pivotal Trials for New
Molecular Entity Drugs.''} \emph{The American Journal of Therapeutics}
22 (6): 435--55.

\bibitem[\citeproctext]{ref-hammer1997controlled}
Hammer, Scott M., Kathleen E. Squires, Michael D. Hughes, Janet M.
Grimes, Lisa M. Demeter, Judith S. Currier, Joseph J. Jr. Eron, et al.
1997. {``A Controlled Trial of Two Nucleoside Analogues Plus Indinavir
in Persons with Human Immunodeficiency Virus Infection and CD4 Cell
Counts of 200 Per Cubic Millimeter or Less.''} \emph{New England Journal
of Medicine} 337 (11): 725--33.
\url{https://doi.org/10.1056/NEJM199709113371101}.

\bibitem[\citeproctext]{ref-Hsieh2000}
Hsieh, Felix Y., and Philip W. Lavori. 2000. {``Sample-Size Calculations
for the Cox Proportional Hazards Regression Model with Nonbinary
Covariates.''} \emph{Controlled Clinical Trials} 21 (6): 552--60.
\url{https://doi.org/10.1016/S0197-2456(00)00104-5}.

\bibitem[\citeproctext]{ref-khan2020women}
Khan, Muhammad Shahzeb, Izza Shahid, Tariq Jamal Siddiqi, Safi U. Khan,
Haider J. Warraich, Stephen J. Greene, Javed Butler, and Erin D. Michos.
2020. {``Ten-Year Trends in Enrollment of Women and Minorities in
Pivotal Trials Supporting Recent US FDA Approval of Novel
Cardiometabolic Drugs.''} \emph{Journal of the American Heart
Association} 9 (10): e015594.
\url{https://doi.org/10.1161/JAHA.119.015594}.

\bibitem[\citeproctext]{ref-moore2016_ch2}
Moore, Dirk F. 2016. {``Basic Concepts and Survival Distributions.''} In
\emph{Applied Survival Analysis Using r}. Springer.
\url{https://doi.org/10.1007/978-3-319-31245-3_2}.

\bibitem[\citeproctext]{ref-poon2013participation}
Poon, Raymond, Khushbu Khanijow, Sabrina Umarjee, Ekaterina Fadiran, M.
Yu, Lei Zhang, and Amish Parekh. 2013. {``Participation of Women and Sex
Analyses in Late-Phase Clinical Trials of New Molecular Entity Drugs and
Biologics.''} \emph{Journal of Women's Health} 22 (7): 604--16.

\bibitem[\citeproctext]{ref-R-base}
R Core Team. 2026. \emph{R: A Language and Environment for Statistical
Computing}. Vienna, Austria: R Foundation for Statistical Computing.
\url{https://www.R-project.org/}.

\bibitem[\citeproctext]{ref-Schoenfeld1983}
Schoenfeld, David A. 1983. {``Sample-Size Formula for the
Proportional-Hazards Regression Model.''} \emph{Biometrics} 39 (2):
499--503. \url{https://doi.org/10.2307/2531021}.

\bibitem[\citeproctext]{ref-GLDreg}
Su, Shengli et al. 2023. \emph{GLDreg: Generalized Lambda Distribution
Regression}. \url{https://CRAN.R-project.org/package=GLDreg}.

\bibitem[\citeproctext]{ref-survival-package}
Therneau, Terry M. 2026. \emph{A Package for Survival Analysis in r}.
\url{https://CRAN.R-project.org/package=survival}.

\end{CSLReferences}




\end{document}
